\section{Discussion}\label{sec:discussion}

\para{A clean dissociation.} Our experiments separate two questions that the phrase ``different kinds of drama'' quietly conflates. First, is the drama \emph{construct} one thing or many? We find it is genuinely multi-dimensional and \emph{shared} across readers: drama decodes perfectly (AUROC $1.000$, \tabref{tab:e1}), yet melodrama is not just more drama---the dramatic$\to$melodramatic step rotates to $83^\circ$ from the drama axis at the final layer (\tabref{tab:e1}), and an eight-subconcept subspace reaches participation ratio $\sim$$5.6$ (vs.\ $\sim$$8.0$ random, vs.\ $1$ for a single axis; \figref{fig:geometry}). Second, do different \emph{people} read along different directions, or do they slide a threshold along the same axis? Here the answer is overwhelmingly the threshold: a per-reader intercept lifts held-out AUROC from $0.660$ to $0.700$ (in-domain, layer 21; \tabref{tab:e3indomain}), while letting each reader choose her own direction in a $30$-dimensional drama subspace adds only $+0.0076$---exactly the shuffled-reader null ($z\!=\!1.0$, n.s.). So ``different kinds of drama'' are \emph{content directions} (earned drama vs.\ melodramatic excess), but ``different people''---for both \goemotions{} intensity disagreement and six prompted personas---are a \emph{sliding threshold} on a fixed, shared axis.

\para{Why this is not trivial.} A reader who suspects drama is multi-dimensional and reader-dependent could be right two ways, and only one survives. The construct could carry an off-axis excess dimension, which it does: \Secref{sec:results} (E1/E2) confirms melodrama as a near-orthogonal excess direction. Or---the stronger H2 claim---individual readers could decode drama from \emph{personal} directions inside that subspace. They do not. The reader-direction model fails its pre-registered control: its gain matches the shuffled-reader permutation null at every layer and scale (layer 21 $+0.0076$ vs.\ $+0.0076$; layer 28 $+0.0046$ vs.\ $+0.0046$; $1.5$B $+0.0084$ vs.\ $+0.0084$), meaning the apparent improvement is pure model capacity that persists even when reader identities are randomized. The multidimensionality lives in the \emph{stimulus} representation and is available to every reader equally; what individuals select is a threshold, not a private axis. Six personas confirm this from the prompt side: their re-derived drama axes are essentially identical (pairwise cosine mean $0.99$, min $\geq0.974$; \figref{fig:persona}), shifting projection and threshold but never the direction.

\para{Relation to prior work.} The shared-axis half of our result sits comfortably with low-dimensional steerable affect manifolds \citep{reichman2025emotions} and with sentiment as essentially one linear direction \citep{tigges2023linear}: a single drama axis decodes near-perfectly and is stable across personas. The reader half refines \citet{orlikowski2025beyond}: reader effects are real and learnable---our $M_0\!\to\!M_1$ threshold lift of $+0.04$ to $+0.11$ is the dominant gain---but they are not captured by low-dimensional demographic structure, and we add that they are not a reader-specific \emph{representational} direction either. The melodrama-as-excess axis is a drama-specific instance of concept-cone and subspace geometry \citep{wollschlager2025geometry}: even an affective construct that decodes as ``one thing'' needs more than one dimension to separate its kinds---just not a reader-conditioned one.

\para{Limitations.} Our verdict is bounded by its operationalizations. For \goemotions{} we operationalize drama as perceived arousal/intensity, and reader labels are binary high-arousal indicators; the H1 conclusion is therefore ``no \emph{detectable} reader-specific direction,'' not proof of absence---graded or continuous scales could expose direction effects the binary cannot resolve. The controlled \dramacorpus{} consists of LLM-authored register variants, and the stimuli are Claude-authored; both were adversarially blind-validated ($290/290$ at $100\%$ agreement), and crucially the subjects of study---Qwen activations and human raters---are independent of the generator, but generator priors could still shape the excess axis. We study one model family (Qwen2.5), replicated across $1.5$B and $7$B and across layers, but cross-family generalization (Llama, Gemma) is untested. Finally, activations are mean-pooled over tokens; token-level probes, SAE decompositions, and causal steering are planned but not run here.

\para{Broader implications.} The practical upshot is that personalizing affect judgments to a reader is often cheap. If individual variation in what counts as ``dramatic'' is mostly a threshold shift on a shared axis rather than steering within a multi-dimensional subspace, then adapting to a reader can be a one-parameter calibration rather than learning a per-user representation---a far smaller and more auditable intervention. At the same time, the construct's own multidimensionality is consequential: earned drama and melodramatic excess occupy nearly orthogonal directions, so moderation, recommendation, and criticism systems that collapse them into a single ``intensity'' score will systematically conflate genuine stakes with sentimentality. The lesson cuts both ways---model the \emph{content} as multi-dimensional, but the \emph{reader} as a movable threshold.
